\documentclass{report}
\usepackage[utf8]{inputenc}
\usepackage[T1]{fontenc}
\usepackage{geometry}
\usepackage{graphicx}
\usepackage{float}
\usepackage{hyperref}
\usepackage[english,indonesian]{babel}

\title{Optimasi Pemrosesan Template LaTeX Menggunakan Pendekatan Berbasis AI}
\author{Academic Support Agent}
\date{\today}

\begin{document}
\maketitle

\begin{abstract}

\end{abstract}

\chapter{PENDAHULUAN}


\chapter{TINJAUAN PUSTAKA}


% Original content (commented out):
% \documentclass[11pt,letterpaper]{report}
\oddsidemargin=-0.2cm %margen izquierdo
\topmargin=0cm   %margen superior
%\topmargin=-1.5cm   %margen superior
%\topmargin=-7.8cm % Cuando se usa 'hyperref'
\textwidth=17cm     %ancho
\textheight=23cm    %alto 
\parindent=0.4cm    %sangria
\parskip=2mm        %espacio entre parrafos
\columnsep=7mm      %espacio entre columnas
\addtolength{\voffset}{-1cm}
\usepackage{stackengine}
\newcommand\mathcalbb[2][2]{%
  \stackengine{0pt}{$\mathcal{#2}$}{$\mkern#1mu\mathcal{#2}$}{O}{l}{F}{F}{L}}
\usepackage{fancyhdr}

%------------------------------------------------------------
%		       				GENERAL PACKAGES
%------------------------------------------------------------
\usepackage[dvips]{graphics}
\usepackage[utf8]{inputenc}
\usepackage[spanish]{babel}
\pagestyle{fancy} % encabezados
\fancyhf{}
\rhead{\thepage}
\usepackage{multicol}
\usepackage{amsmath,amsthm,amssymb,psfrag,epsfig,float,amsfonts,latexsym,color}
\usepackage[ruled]{algorithm2e}
\usepackage{booktabs}
\usepackage{graphicx}
\usepackage{tabu}
\usepackage{bbm}
\usepackage{multirow}	 % uniendo varias filas
\usepackage{arydshln}	 % dashed line in a table
\usepackage{indentfirst} % poner sangria en el primer parrafo del reporte
\usepackage{mathrsfs}    % \mathscr
\usepackage[T1]{fontenc}
\usepackage{ifthen}
\usepackage{enumerate}
\usepackage{longtable}   % tablas largas que se pasan de una hoja
\usepackage{siunitx}
\usepackage[utf8]{inputenc}

\usepackage{fancybox}
\usepackage{stmaryrd}
\usepackage{url}
\usepackage[table]{xcolor}
\usepackage{longfbox}
\newcommand{\urria}[1]{\lfbox[border-style={none,solid,solid,none}]{#1}}
%\usepackage{refcheck}  % Paquete para revisar las referencias
%------------------------------------------------------------
%		LINKS (Agregar Hyperlinks)
%------------------------------------------------------------
\def\phantomsection{}
\usepackage{url}
\usepackage[colorlinks=true,citecolor=red,linkcolor=black,
urlcolor=black,bookmarks=true,breaklinks=true]{hyperref}
%------------------------------------------------------------
%			DEFINICIONES
%------------------------------------------------------------
% Para los encabezados
\renewcommand{\qedsymbol}{$\blacksquare$}
\renewcommand{\baselinestretch}{1.1}
\renewcommand{\chaptermark}[1]{%
	\markboth{{\small\textsf{\chaptername
				\ \thechapter.\ #1}}}{}}
\renewcommand{\sectionmark}[1]{\markright{\small\textsf{\thesection.\ #1}}}%
% Nombre de Secciones
\def\bibname{References}
% Para el estilo de los capitulos
\usepackage{titlesec} 
\titleformat{\chapter}[display]
{\bfseries\Large}
{\filright\MakeUppercase{\chaptertitlename} \Huge\thechapter}
{1ex}
{\titlerule\vspace{1ex}\filleft}
[\vspace{1ex}\titlerule]

\renewcommand{\headrulewidth}{0pt}
\newcommand{\cred}{}
\usepackage{etoolbox}



\makeatletter
\patchcmd{\ttlh@hang}{\parindent\z@}{\parindent\z@\leavevmode}{}{}
\patchcmd{\ttlh@hang}{\noindent}{}{}{}
\makeatother

%============================================================

%%%%%%%%%  OUR DEFINITIONS %%%%%%%%%%%%
\newtheorem{thm}{Theorem}[chapter]
\newtheorem{rem}{Remark}[chapter]
\newtheorem{lemma}[thm]{Lema}
\newtheorem{cor}[thm]{Corolario}
\newtheorem{defi}{Definición}[chapter]
\newtheorem{axio}{Axioma ZF}
\newtheorem{teo}{Teorema}[chapter]
\newtheorem{assump}[thm]{Assumption}
\newtheorem{rmrk}[thm]{Remark}



\newcommand{\ds}{\displaystyle}
\newcommand{\alphan}{{\boldsymbol{\alpha}}}

\newcommand{\proofend}{\hfill$\square$}
\newcommand\disp{\displaystyle}
\newcommand{\punto}{\hspace{2pt}\cdot\hspace{2pt}}
\newcommand{\ts}{\textstyle}
\newcommand{\smallfrac}[2]{{\textstyle\frac{#1}{#2}}}

\def\qin{{\quad\hbox{en}\quad}}
\def\qon{{\quad\hbox{sobre}\quad}}
\def\qan{{\quad\hbox{y}\quad}}
\def\inn#1{\langle{#1}\rangle}
\newcommand{\pin}{\mathbf{\Pi}}
%%%%%%% GREEK LETTERS %%%%%%%

\newcommand{\Om}{\Omega}
\newcommand{\g}{{\gamma}}
\newcommand{\G}{\Gamma}
\newcommand{\kh}{\hat{K}}
\newcommand{\bnu}{{\boldsymbol{\nu}}}
\newcommand{\bphi}{{\boldsymbol\phi}}
\newcommand{\bvarphi}{{\mbox{\boldmath $\varphi$}}}
\newcommand{\nmn}{{\boldsymbol{n}}}
\newcommand{\taumn}{\btau}
\newcommand{\bva}{{\boldsymbol\varphi}}
\newcommand{\bcero}{{\boldsymbol{0}}}
\newcommand{\zmn}{\bzeta}
%%%%%%%%%% MATHBF %%%%%%%%%%%%%%%%%%%%
\def\bA{\mathbf{A}}
\def\bC{\mathbf{C}}
\def\bF{\mathbf{F}}
\def\bK{\mathbf{K}}
\def\bI{\mathbf{I}}
\def\bX{\mathbf{X}}
\def\bP{\mathbf{P}}
\def\bV{\mathbf{V}}
\def\bM{\mathbf{M}}
\def\bR{\mathbf{R}}
\def\bS{\mathbf{S}}
\def\bT{\mathbf{T}}
%%%%%%%%%%%%%%% MATHRM %%%%%%%%%%%%%%%%%%%%%%
\newcommand\rmF{\mathrm{F}}
\newcommand\rmH{\mathrm{H}}
\newcommand\rmQ{\mathrm{Q}}
\newcommand{\R}{\mathrm{R}}
\newcommand{\Prm}{\mathrm{P}}
\newcommand{\rmL}{\mathrm{L}}
\newcommand{\rmC}{\mathrm{C}}
%%%%%%%%%%%%%%% MATHBB %%%%%%%%%%%%%%%%%Conclusions%%%%%
\newcommand{\bbK}{{\mathbb{K}}}
\newcommand\bbM{\mathbb{M}}
\newcommand\bbI{\mathbb{I}}
\newcommand\Ibb{\mathbb{I}}
\newcommand\bbP{\mathbb{P}}
\newcommand\bbH{\mathbb{H}}
\newcommand\bbL{\mathbb{L}}
\newcommand{\bbA}{\mathbb{A}}

%%%%%%%%%%%%%%% MATHCAL %%%%%%%%%%%%%%%%%%%%%%
\newcommand{\Qcal}{{\mathcal{Q}}}
\newcommand{\Bcal}{{\mathcal{B}}}
\newcommand{\Bcaln}{{\boldsymbol{\Bcal}}}
\newcommand{\Gcaln}{{\boldsymbol{\Gcal}}}
\newcommand{\Ocal}{{\mathcal{O}}}
\newcommand{\Pcal}{{\mathcal{P}}}
\newcommand{\Pcaln}{{\boldsymbol{\Pcal}}}
\newcommand{\Scal}{\mathcal{S}}
\newcommand{\Tcal}{\mathcal{T}}
\newcommand{\Ecal}{\mathcal{E}}
\newcommand{\Pcalbb}{\mathcalbb{P}}
\newcommand{\Rcaln}{{\boldsymbol{\Rcal}}}
%%%% greek bold letters
\newcommand{\bgamma}{{\boldsymbol\gamma}}
\newcommand{\bkappa}{{\boldsymbol\kappa}}
\newcommand{\blambda}{{\boldsymbol\lambda}}
\newcommand{\bLambda}{{\boldsymbol\Lambda}}
\newcommand{\bbeta}{{\boldsymbol\eta}}
\newcommand{\bsi}{{\boldsymbol\sigma}}
\newcommand{\bPsi}{\boldsymbol{\Psi}}
\newcommand{\bPhi}{{\boldsymbol\Phi}}
\newcommand{\bpsi}{{\mbox{\boldmath $\psi$}}}
\newcommand{\bPi}{{\mbox{\boldmath $\Pi$}}}
\newcommand{\btau}{{\boldsymbol\tau}}
\newcommand{\bps}{{\boldsymbol\psi}}
\newcommand{\bchi}{{\boldsymbol\chi}}
\newcommand{\bxi}{{\boldsymbol\xi}}
\newcommand{\brho}{{\boldsymbol\rho}}
\newcommand{\bzeta}{{\boldsymbol\zeta}}
%%%% bold letters
\renewcommand{\b}{\mathbf{b}}
\newcommand{\ubsi}{\underline{\bsi}}
\newcommand{\ubtau}{\underline{\btau}}
\newcommand{\ubphi}{\underline{\bphi}}
\newcommand{\ubpsi}{\underline{\bpsi}}
\newcommand{\ubzeta}{\underline{\bzeta}}

%------------------------------------------------------------

\begin{document}
%============================================================
%						PRELIMINARIES
%============================================================

\include{preliminaries/front_cover}

\renewcommand{\thepage}{\roman{page}} 
\setcounter{page}{2}

\include{preliminaries/resumen}

\include{preliminaries/abstract}

\include{preliminaries/agradecimientos}


%============================================================
%                        	INDEX
%============================================================
% Tabla de contenidos, lista de tablas y figuras
\clearpage
%--------------------------------------------------- hyperref
\phantomsection
\addcontentsline{toc}{chapter}{Contenidos}
\tableofcontents
%--------------------------------------------------- hyperref
\clearpage

%============================================================
%                 FIN DE LA NUMERACION ROMANA
%============================================================
\clearpage
\renewcommand{\thepage}{\arabic{page}} % Numeracion arabica
\setcounter{page}{1}

%============================================================
%                 CHAPTERS OF THE THESIS
%============================================================
\setcounter{chapter}{-1}

\include{chapters/introduccion}
\include{chapters/ch-0}
\include{chapters/ch-1}
\include{chapters/ch-2}
\include{chapters/ch-3}
\include{chapters/ch-4}
%\phantomsection
%\addcontentsline{toc}{chapter}{Conclusions}

\pagestyle{fancy}
\fancyhf{}
%\fancyhead[RE,LO]{\small\textsf{Conclusions and Future Works}}
\fancyhead[LE,RO]{\thepage}



%\phantomsection
%\addcontentsline{toc}{chapter}{Conclusiones}

%\pagestyle{fancy}
%\fancyhf{}
%\fancyhead[RE,LO]{\small\textsf{Conclusiones y Trabajos Futuros}}
%\fancyhead[LE,RO]{\thepage}

\include{chapters/conclusiones_trabajos-futuros}

\clearpage
%--------------------------------------------------- hyperref
\phantomsection
\addcontentsline{toc}{chapter}{Referencias}

%\pagestyle{fancy}
%\fancyhf{}
%\fancyhead[RE,LO]{\small\textsf{References}}
%\fancyhead[LE,RO]{\thepage}

%------------------------------------------------------------
\nocite{*} % Aparecen todas las referencias aunque no se citen
\bibliography{chapters/references}
\bibliographystyle{unsrt}
\include{chapters/references}
\include{chapters/glosario}
\end{document}

\end{document}